\section{Expanded Cross-Fitted Attacks and Oracle Localization}\label{sec:r7-crossfit-supp}
This extension reanalyzes the existing 690 Sol episodes with the same five customer folds as E10. It is post hoc: the cohort, defense design and prior results were already available. The new source/configuration freeze precedes the new fits, not the original test collection. All 53 customers are retained, with 14 cancellation-positive and 18 return-positive customers. Each prompt contributes 345 out-of-fold episodes. Fold sizes are 11/12/10/10/10 customers; a rare joint-label stratum has only three customers, so not every fold can contain that stratum. The saved fold assignment and customer/target counts make this restriction explicit.

\subsection{Complete attack union and fitting boundary}
For each fold, joint count templates, cover donors, framing caps and size-to-tool decoders use only the other four folds. The extended union retains all nine prior network attacks: majority, undefended-data-trained gradient boosting, D2 logistic regression and gradient boosting, random forest, nearest neighbors, and three decoded-profile attacks. It adds the fixed trained 1D-CNN and k-FP-style leaf attack described in S20. CNN inputs are direction, log record length and a mask; the timing channel is identically zero because this analysis concerns length/order. No interarrival times are synthesized. The mapping from archived C/S direction strings to numeric signs is explicit and checked. An initial mapping error stopped before any prediction pack was produced; its failed protocol snapshot is retained separately.

All 60 fold/prompt/condition/target packs preserve their fitted learners, test scores and predictions. Out-of-fold rows are pooled only after every fold is evaluated. Each of 2,000 common customer-bootstrap draws retains a customer's tasks and repetitions together and reselects the full 11-attack maximum. Paired decreases use the same draw for undefended and defended conditions. These intervals condition on the fitted attackers and original cohort. The absolute gate remains an upper BA bound below .60 and suite-maximum AUROC below .65. Neither an attack's favorable point estimate nor an increase in the number of resampled trajectories replaces independent positive customers.

\begin{longtable}{@{}lllrrrr@{}}
\caption{Expanded E10: all 12 cells. Percent BA and paired undefended-minus-defended decrease, each with 95\% customer-bootstrap intervals. AUROC is independently maximized over the same 11 attacks. Ratios are bytes/calls.}\label{tab:r7-crossfit-all}\\
\toprule Prompt & Condition & Target & BA [95\% CI] & AUROC & Decrease [95\% CI] & B$\times$/C$\times$\\\midrule\endfirsthead
\toprule Prompt & Condition & Target & BA [95\% CI] & AUROC & Decrease [95\% CI] & B$\times$/C$\times$\\\midrule\endhead
Standard & Undefended & cancel & 91.7 [82.1, 100.0] & 0.979 & -- & 1.00/1.00\\
Standard & Undefended & return & 85.5 [76.7, 93.4] & 0.897 & -- & 1.00/1.00\\
Standard & Joint constant & cancel & 60.5 [52.9, 72.4] & 0.594 & 31.2 [15.3, 41.0] & 19.76/5.04\\
Standard & Joint constant & return & 55.5 [51.8, 64.2] & 0.533 & 30.0 [14.8, 39.2] & 19.76/5.04\\
Standard & Joint buckets & cancel & 57.1 [50.0, 71.5] & 0.579 & 34.6 [15.1, 46.4] & 12.36/5.04\\
Standard & Joint buckets & return & 52.7 [50.0, 61.7] & 0.536 & 32.8 [20.8, 38.9] & 12.36/5.04\\
Rules & Undefended & cancel & 90.7 [81.5, 98.9] & 0.978 & -- & 1.00/1.00\\
Rules & Undefended & return & 85.5 [76.6, 93.6] & 0.880 & -- & 1.00/1.00\\
Rules & Joint constant & cancel & 58.8 [52.1, 71.6] & 0.581 & 32.0 [13.1, 44.2] & 19.88/5.04\\
Rules & Joint constant & return & 59.4 [54.2, 68.8] & 0.581 & 26.1 [11.1, 36.6] & 19.88/5.04\\
Rules & Joint buckets & cancel & 56.9 [50.2, 71.2] & 0.573 & 33.8 [14.2, 45.8] & 12.45/5.04\\
Rules & Joint buckets & return & 52.7 [50.0, 65.0] & 0.565 & 32.8 [19.0, 39.8] & 12.45/5.04\\
\bottomrule\end{longtable}


\subsection{What changes under stronger attacks}
In the standard-prompt arm, the CNN raises bucketed cancellation inference to 57.1\%, and the leaf-fingerprint learner raises constant-framing return inference to 55.5\%. The corresponding paired decreases from undefended transport remain positive: 34.6 points [15.1,46.4] and 30.0 points [14.8,39.2]. This exploratory comparison shows measured reduction with an unresolved absolute bound. It is not tied to the original 50\% frozen-split realization. Both prompt arms and every evaluated operating point are retained, including rules-arm estimates that differ from the standard prompt. The main descriptive table reports the complete union; S15 retains the earlier suite and six-point selection analysis as historical context.

The extension saves 4,140 out-of-fold condition/target rows and checks all 45,540 attack scores against the fitted model archives. Input hashes preserve the original cohort and folds; a documented path map permits read-only verification after checkout relocation without changing frozen source paths or data hashes.

\subsection{Expanded oracle localization}
A separate before-fit extension reuses the original E12 oracle filter and the same five folds. It retains or removes complete exchanges belonging to true state-changing calls, then exposes only network record directions and lengths to attackers. Tool identities choose the diagnostic subset but are not attack features. The full-view reference is the verified expanded E10 undefended output. All nine prior attacks plus the same fixed CNN and forest-leaf learners are fitted for each oracle view, yielding 40 additional fold packs and 2,760 out-of-fold rows. Exact model reloading and legacy refitting check 30,360 scores, customer separation, labels and the full-view reference. Empty filtered traces remain included.

\begin{longtable}{@{}lllrr@{}}
\caption{Expanded E12 oracle ablation, both prompts. Suite-maximum BA (percent) and AUROC, with customer-bootstrap 95\% intervals. The full reference is the verified expanded E10 undefended prediction set.}\label{tab:r7-e12-all}\\
\toprule Prompt & Target & Oracle view & BA [95\% CI] & AUROC [95\% CI]\\\midrule\endfirsthead
\toprule Prompt & Target & Oracle view & BA [95\% CI] & AUROC [95\% CI]\\\midrule\endhead
Standard & cancel & Full trace & 91.7 [82.1, 100.0] & 0.979 [0.949, 1.000]\\
Standard & cancel & Writes removed & 62.8 [54.6, 75.3] & 0.611 [0.533, 0.768]\\
Standard & cancel & Writes only & 91.7 [82.1, 100.0] & 0.973 [0.936, 1.000]\\
Standard & return & Full trace & 85.5 [76.7, 93.4] & 0.897 [0.849, 0.965]\\
Standard & return & Writes removed & 59.9 [52.2, 71.8] & 0.657 [0.571, 0.789]\\
Standard & return & Writes only & 84.9 [78.8, 94.0] & 0.906 [0.864, 0.973]\\
Rules & cancel & Full trace & 90.7 [81.5, 98.9] & 0.978 [0.954, 0.999]\\
Rules & cancel & Writes removed & 68.7 [59.4, 80.2] & 0.690 [0.593, 0.822]\\
Rules & cancel & Writes only & 90.7 [81.4, 98.9] & 0.968 [0.930, 0.998]\\
Rules & return & Full trace & 85.5 [76.6, 93.6] & 0.880 [0.805, 0.955]\\
Rules & return & Writes removed & 66.5 [56.5, 77.7] & 0.664 [0.579, 0.807]\\
Rules & return & Writes only & 85.1 [78.6, 93.7] & 0.917 [0.862, 0.978]\\
\bottomrule\end{longtable}


For the standard prompt, write exchanges constitute 19.6\% of real calls and 23.4\% of TLS wire bytes, yet retain 91.7/84.9\% BA versus 91.7/85.5\% for full traces. Removing writes reduces the suite maximum by 28.9 points [10.7,41.9] for cancellation and 25.6 [11.0,35.4] for returns. The rules-arm decreases are 22.0 [6.0,35.1] and 19.0 [6.0,30.1] points. Read-only traces still carry measurable information; this is concentration of the main signal, not exclusive localization. The small full-minus-writes-only point difference does not constitute a formal equivalence test. Per-prompt read/write retained-call and byte fractions sum to one and are checked against the original trace projection.
